\documentclass[11pt]{article}
\usepackage{graphicx}
\usepackage{amssymb}
\usepackage{bm}
\usepackage{mathtools}
\usepackage{amsmath}
\usepackage{commath}
\usepackage{amsthm}
\usepackage{enumitem}
\usepackage{float} 
\usepackage{array}
\usepackage{outlines}
\usepackage{setspace}
\usepackage[colorinlistoftodos]{todonotes}
\setlist[enumerate,1]{label=\alph*)}
\setlist[enumerate,2]{label=\roman*)}
\newtheorem{theorem}{Theorem}
\newtheorem{lemma}{Lemma}
\newtheorem{remark}{Remark}
\newtheorem{definition}{Definition}
\newtheorem{prop}{Proposition}
\newcommand*\diff{\mathop{}\!\mathrm{d}}
\usepackage{tikz}
\usetikzlibrary{arrows,shapes,trees, calc}

%%% *** PREAMBLE *** %%%
\title{Things to consider, ask, include and look into.}

%%%***  CONTENT *** %%%
\begin{document}

\maketitle

\begin{itemize}

\item \textbf{How to interpret the "3D"-ness of the limit model with $\epsilon = 0:$} it represents that the vertical dimension is much smaller, compared to the horizontal dimensions. It is a relative thing.

\item \textbf{The divergence-free condition and the Aubin-Lions lemma.} In the A-G article it is mentioned that the Aubin-Lions lemma can not be used because the velocity including $\epsilon$ is not divergence-free. Actually, being divergence-free is not a requirement, it is simply easier when we have it, and without it the conditions of the original A-L lemma can not be verified. This is why they use another version of it, but not because being divergence-free is a condition in itself.

\item \textbf{Why do we have $\epsilon^3$ in the Gaussian:} Because it is a $3D$ Gaussian. The number of dimensions appears there.

\item \textbf{Weak formulation of $\mathbf{u}^\epsilon$}: $\mathbf{u}^\epsilon$ is defined to be in $L^2 H^1$ because at this point we do all this for a fix $\epsilon,$ which is as if we simply had a constant. On the other hand, it makes sense to define it like this, because in this case the system we have is essentially the complete Navier-Stokes, for which we have the results that verify the existence of such a weak solution.

\item \textbf{The structure of stating Proposition 1.:} Let $\mathbf{u}^\epsilon$ be a weak solution of the anisotropic system, then we state results such as for example $\mathbf{u}^\epsilon \in L^2 L^2.$ In this case we have switched from a fix $\epsilon$ to an $\epsilon$ tending to zero, so this is already a different scenario, that is why it is not a contradiction to use the definition of WS (which, for a fix epsilon, had $H^1$ regularity), and then state $L^2 L^2.$
 
\end{itemize}

\end{document}

%%% add http://elte.prompt.hu/sites/default/files/tananyagok/AtmosphericChemistry/ch10s03.html
%%% http://twister.caps.ou.edu/PM2000/Chapter7.pdf
